%% ============================================================
%%  Detailed Calculation Breakdown: Topological Magnons in Y2V2O7
%%  Companion Document for Advisor Presentation
%%  Author: Liam T. Schmidt
%% ============================================================
\documentclass[11pt,letterpaper]{article}

%% --- Packages ---
\usepackage[margin=1in]{geometry}
\usepackage{amsmath,amssymb,amsfonts}
\usepackage{bm}
\usepackage{graphicx}
\usepackage[colorlinks=true,linkcolor=blue,citecolor=blue,urlcolor=blue]{hyperref}
\usepackage{siunitx}
\usepackage{physics}
\usepackage{xcolor}
\usepackage{booktabs}
\usepackage{enumitem}
\usepackage{tcolorbox}
\usepackage{fancyhdr}

\tcbuselibrary{theorems,skins,breakable}

%% Custom environments for key results and caveats
\newtcolorbox{keyresult}[1][]{
  colback=green!5!white,
  colframe=green!50!black,
  title={Key Result},
  fonttitle=\bfseries,
  #1
}

\newtcolorbox{assumption}[1][]{
  colback=yellow!10!white,
  colframe=orange!70!black,
  title={Assumption / Caveat},
  fonttitle=\bfseries,
  #1
}

\newtcolorbox{derivation}[1][]{
  colback=blue!5!white,
  colframe=blue!50!black,
  title={Derivation Details},
  fonttitle=\bfseries,
  breakable,
  #1
}

\pagestyle{fancy}
\fancyhf{}
\rhead{Topological Magnons Calculations}
\lhead{Liam Schmidt}
\cfoot{\thepage}

\sisetup{separate-uncertainty=true, range-phrase={--}, range-units=single}

%% ============================================================
\begin{document}

\title{\textbf{Detailed Calculation Breakdown}\\[0.5em]
\Large Topological Magnon Band Structure in Y$_2$V$_2$O$_7$:\\
Weyl Magnons and Polarimetric RIXS Signatures}

\author{Liam T. Schmidt\\
\small Department of Physics, Northeastern University}

\date{\today}

\maketitle

\tableofcontents
\newpage

%% ============================================================
\section{Overview and Scope}
%% ============================================================

This document provides a detailed breakdown of all theoretical calculations presented in the manuscript ``Topological Magnon Band Structure in Y$_2$V$_2$O$_7$: Weyl Magnons and Polarimetric RIXS Signatures.'' Each calculation is explained step-by-step, with explicit attention to:

\begin{itemize}
    \item Physical assumptions and their justifications
    \item Mathematical derivations and approximations
    \item Numerical implementation details
    \item Parameter choices and their sources
    \item Potential limitations and caveats
\end{itemize}

The calculations fall into four main categories:
\begin{enumerate}
    \item \textbf{Linear Spin-Wave Theory (LSWT)} --- Magnon band structure
    \item \textbf{Topological Analysis} --- Berry curvature and Chern numbers
    \item \textbf{$k \cdot p$ Theory} --- Weyl point characterization
    \item \textbf{Cluster Exact Diagonalization (CED)} --- RIXS spectra
\end{enumerate}

%% ============================================================
\section{Spin Hamiltonian Construction}
%% ============================================================

\subsection{The Full Hamiltonian}

The starting point is the spin Hamiltonian for nearest-neighbor V$^{4+}$ ions on the pyrochlore lattice:

\begin{equation}
\boxed{
\mathcal{H} = -J \sum_{\langle ij \rangle} \bm{S}_i \cdot \bm{S}_j
             + \sum_{\langle ij \rangle} \bm{D}_{ij} \cdot
               \left(\bm{S}_i \times \bm{S}_j\right)
}
\label{eq:Hfull}
\end{equation}

\begin{derivation}[title={Term-by-term breakdown}]
\textbf{Heisenberg Exchange (First Term):}
\begin{equation}
\mathcal{H}_{\text{Heis}} = -J \sum_{\langle ij \rangle} \bm{S}_i \cdot \bm{S}_j
= -J \sum_{\langle ij \rangle} \left( S_i^x S_j^x + S_i^y S_j^y + S_i^z S_j^z \right)
\end{equation}

The sign convention: $J > 0$ means \emph{ferromagnetic} coupling (parallel spins lower energy). This is opposite to the antiferromagnetic convention where $J > 0$ favors antiparallel alignment.

\textbf{Dzyaloshinskii-Moriya (DM) Interaction (Second Term):}
\begin{equation}
\mathcal{H}_{\text{DM}} = \sum_{\langle ij \rangle} \bm{D}_{ij} \cdot \left(\bm{S}_i \times \bm{S}_j\right)
\end{equation}

Expanding the cross product:
\begin{equation}
\bm{S}_i \times \bm{S}_j = 
\begin{pmatrix}
S_i^y S_j^z - S_i^z S_j^y \\
S_i^z S_j^x - S_i^x S_j^z \\
S_i^x S_j^y - S_i^y S_j^x
\end{pmatrix}
\end{equation}

The DM interaction is \textbf{antisymmetric}: $\bm{D}_{ij} = -\bm{D}_{ji}$, so the ordering of sites $i,j$ matters.
\end{derivation}

\subsection{DM Vector Constraints from Symmetry}

\begin{assumption}[title={Moriya's Rules Applied to Pyrochlore}]
In the pyrochlore structure ($Fd\bar{3}m$, No.~227):
\begin{itemize}
    \item Each nearest-neighbor bond lies on a $C_2$ axis
    \item The bond midpoint is the inversion center for that axis
    \item $\bm{D}_{ij}$ must be perpendicular to the bond direction
    \item $\bm{D}_{ij}$ lies in the plane containing the bond and the $C_2$ axis
\end{itemize}

The six DM vectors on a single tetrahedron are related by $C_3$ rotation about $[111]$. Only the magnitude $D = |\bm{D}_{ij}|$ is a free parameter.
\end{assumption}

\subsection{Parameter Values and Sources}

\begin{table}[h]
\centering
\begin{tabular}{lccc}
\toprule
Parameter & Value & Units & Source \\
\midrule
$J$ (exchange) & 8.22 & meV & Onose et al. 2010 (INS on Lu$_2$V$_2$O$_7$) \\
$D/J$ (DM ratio) & 0.32 & -- & Thermal Hall fit (Lu$_2$V$_2$O$_7$) \\
$D$ (DM strength) & 2.63 & meV & $D = 0.32 \times J$ \\
$a$ (lattice constant) & 9.89 & \AA & Gardner et al. 2010 \\
$T_C$ (Curie temp.) & 70 & K & Ideue et al. 2012 \\
\bottomrule
\end{tabular}
\caption{Model parameters and their experimental sources.}
\label{tab:params}
\end{table}

\begin{assumption}[title={Parameter Transfer Assumption}]
\textbf{Critical caveat:} No direct measurements of $J$ or $D$ exist for Y$_2$V$_2$O$_7$ itself. Parameters are \emph{transferred} from Lu$_2$V$_2$O$_7$ based on:
\begin{itemize}
    \item Similar $T_C \approx 70$ K
    \item Same V$^{4+}$ ($3d^1$, $S=1/2$) electronic configuration
    \item Similar lattice constants (9.89 vs 9.93 \AA)
\end{itemize}

\textbf{Uncertainty:} The V--O--V bond angle differs by $1$--$2^\circ$ due to different A-site ionic radii (Y$^{3+}$ vs Lu$^{3+}$), potentially shifting $D/J$ by $\pm 20\%$.

\textbf{Robustness:} The Weyl crossing persists over $0.1 \lesssim D/J \lesssim 0.6$, so the topology is robust to this uncertainty.
\end{assumption}

%% ============================================================
\section{Holstein-Primakoff Transformation}
%% ============================================================

\subsection{The Transformation}

To map the spin problem to a bosonic magnon problem, we use the Holstein-Primakoff (HP) transformation:

\begin{derivation}[title={Holstein-Primakoff Mapping}]
For a spin $S$ at site $i$, define bosonic operators $a_i, a_i^\dagger$ satisfying $[a_i, a_j^\dagger] = \delta_{ij}$. The exact transformation is:

\begin{align}
S_i^+ &= \sqrt{2S} \sqrt{1 - \frac{a_i^\dagger a_i}{2S}} \, a_i \\
S_i^- &= \sqrt{2S} \, a_i^\dagger \sqrt{1 - \frac{a_i^\dagger a_i}{2S}} \\
S_i^z &= S - a_i^\dagger a_i
\end{align}

where $S^\pm = S^x \pm i S^y$.

\textbf{Linear Spin-Wave Approximation:} Expand the square root to leading order in $1/S$:
\begin{equation}
\sqrt{1 - \frac{n_i}{2S}} \approx 1 - \frac{n_i}{4S} - \frac{n_i^2}{32S^2} - \ldots
\end{equation}

Keeping only the leading term ($n_i/S \ll 1$, i.e., low magnon density):
\begin{equation}
\boxed{
S_i^+ \approx \sqrt{2S}\, a_i, \quad
S_i^- \approx \sqrt{2S}\, a_i^\dagger, \quad
S_i^z = S - a_i^\dagger a_i
}
\end{equation}

This is \textbf{exact} for $S \to \infty$ and an approximation for finite $S$.
\end{derivation}

\begin{assumption}[title={Validity for $S = 1/2$}]
\textbf{Critical caveat:} V$^{4+}$ has $S = 1/2$, the \emph{smallest possible} spin. LSWT corrections scale as $1/S$, so quantum corrections are expected to be $\sim 100\%$ nominally.

\textbf{Mitigating factors:}
\begin{itemize}
    \item Prior $1/S$ studies of Lu$_2$V$_2$O$_7$ show actual renormalization is $\sim 20\%$ due to partial cancellation
    \item The Weyl crossing is \emph{topologically protected} by $C_3$ symmetry --- interactions can shift $t_W$ and $\omega_W$ but \emph{cannot gap} the crossing
\end{itemize}
\end{assumption}

\subsection{Substitution into the Hamiltonian}

\begin{derivation}[title={LSWT Hamiltonian Construction}]
Substituting the HP transformation into $\mathcal{H}$:

\textbf{Heisenberg term:}
\begin{align}
\bm{S}_i \cdot \bm{S}_j &= S_i^z S_j^z + \frac{1}{2}(S_i^+ S_j^- + S_i^- S_j^+) \\
&= (S - n_i)(S - n_j) + S(a_i a_j^\dagger + a_i^\dagger a_j) \\
&\approx S^2 - S(n_i + n_j) + S(a_i a_j^\dagger + a_i^\dagger a_j) + O(a^4)
\end{align}

The constant $S^2$ is the classical ground state energy. The linear terms vanish for the ferromagnetic ground state (no sublattice canting). The bilinear terms give the magnon hopping.

\textbf{DM term:} Similarly generates off-diagonal magnon hopping with complex (imaginary) matrix elements.

The full bilinear Hamiltonian in $k$-space is:
\begin{equation}
\mathcal{H}_{\text{SW}} = \frac{1}{2} \sum_{\bm{k}} \bm{\Psi}_{\bm{k}}^\dagger 
\begin{pmatrix}
A(\bm{k}) & B(\bm{k}) \\
B^*(-\bm{k}) & A^T(-\bm{k})
\end{pmatrix}
\bm{\Psi}_{\bm{k}}
\end{equation}

where $\bm{\Psi}_{\bm{k}}^\dagger = (a_1^\dagger(\bm{k}), \ldots, a_4^\dagger(\bm{k}), a_1(-\bm{k}), \ldots, a_4(-\bm{k}))$ is the Nambu spinor for 4 sublattices.
\end{derivation}

\begin{keyresult}[title={Simplification for Collinear Ferromagnet}]
For the collinear ferromagnetic ground state (all spins parallel along a global axis):
\begin{itemize}
    \item The anomalous terms vanish: $B(\bm{k}) = 0$
    \item The $8 \times 8$ BdG problem reduces to a $4 \times 4$ Hermitian eigenvalue problem
\end{itemize}

This simplification occurs because HP expansion in a uniform frame produces only number-conserving hopping.
\end{keyresult}

%% ============================================================
\section{Bogoliubov Diagonalization}
%% ============================================================

\subsection{The Colpa Algorithm}

\begin{derivation}[title={Paraunitary Diagonalization}]
The BdG Hamiltonian must be diagonalized while preserving the bosonic commutation relations. This requires a \emph{paraunitary} transformation $\mathcal{T}(\bm{k})$ satisfying:
\begin{equation}
\mathcal{T}^\dagger \Sigma_z \mathcal{T} = \Sigma_z
\end{equation}
where $\Sigma_z = \text{diag}(1,1,1,1,-1,-1,-1,-1)$ is the bosonic metric.

\textbf{Colpa's Algorithm (1978):}
\begin{enumerate}
    \item Form $\Sigma_z H(\bm{k})$ where $H$ is the BdG matrix
    \item Compute eigenvalues and eigenvectors of $\Sigma_z H$
    \item Eigenvalues come in $\pm$ pairs; positive ones are magnon energies
    \item Normalize eigenvectors with respect to the $\Sigma_z$ metric
\end{enumerate}

For our collinear FM case with $B = 0$, this reduces to standard Hermitian diagonalization of the $4 \times 4$ matrix $A(\bm{k})$.
\end{derivation}

\subsection{Numerical Implementation Details}

\begin{itemize}
    \item High-symmetry path: $\Gamma \to X \to W \to L \to \Gamma$ (FCC convention)
    \item Sampling: 300 points per segment
    \item Weyl point search: 5000 points along $\Gamma \to L$, gap threshold $0.05$ meV
    \item Magnon Hall conductivity: $20 \times 20 \times 20$ Monkhorst-Pack grid
\end{itemize}

%% ============================================================
\section{Berry Curvature Calculation}
%% ============================================================

\subsection{Definition and Physical Meaning}

The Berry curvature is the ``magnetic field'' in $k$-space:

\begin{equation}
\bm{\Omega}_n(\bm{k}) = i \langle \nabla_{\bm{k}} u_{n\bm{k}} | \times | \nabla_{\bm{k}} u_{n\bm{k}} \rangle
\end{equation}

\begin{derivation}[title={Kubo Formula for Berry Curvature}]
A gauge-invariant expression uses the velocity matrix elements:
\begin{equation}
\boxed{
\Omega_n^{\alpha\beta}(\bm{k}) = -2 \, \text{Im} \sum_{m \neq n}
\frac{\langle n\bm{k} | \partial_{k_\alpha} \mathcal{H} | m\bm{k} \rangle
      \langle m\bm{k} | \partial_{k_\beta} \mathcal{H} | n\bm{k} \rangle}
     {[E_m(\bm{k}) - E_n(\bm{k})]^2}
}
\end{equation}

\textbf{Why gauge-invariant?} Depends only on:
\begin{itemize}
    \item Energy differences $E_m - E_n$ (gauge-invariant)
    \item Off-diagonal velocity matrix elements (invariant under phase changes of eigenstates)
\end{itemize}

\textbf{Numerical implementation:}
\begin{itemize}
    \item Derivatives $\partial_{k_\alpha} \mathcal{H}$ computed by finite differences
    \item Step size: $\delta k = 2 \times 10^{-5}$ \AA$^{-1}$
    \item Sum over $m$ includes all three other magnon branches
\end{itemize}
\end{derivation}

\subsection{Berry Curvature Near a Weyl Point}

\begin{keyresult}[title={Monopole Structure}]
A Weyl point acts as a \emph{Berry monopole} in $k$-space:
\begin{equation}
\bm{\Omega}_n(\bm{k}) \propto \frac{\hat{\bm{r}}}{|\bm{k} - \bm{k}_W|^2}
\end{equation}

This $1/r^2$ divergence is the signature of a topological defect. The color map in Fig.~7a of the paper shows this monopole structure.
\end{keyresult}

%% ============================================================
\section{Chern Number and Topological Charge}
%% ============================================================

\subsection{Definition}

\begin{equation}
\boxed{
\mathcal{C} = \frac{1}{2\pi} \oint_{S^2} \bm{\Omega}_n(\bm{k}) \cdot d\bm{S}
}
\end{equation}

This counts the net Berry flux through any closed surface enclosing the Weyl point.

\begin{derivation}[title={Numerical Integration}]
\textbf{Method:}
\begin{enumerate}
    \item Define spherical coordinates $(\theta, \phi)$ centered at $\bm{k}_W$
    \item Sample $\bm{\Omega}_n \cdot \hat{r}$ on a $20 \times 40$ angular grid
    \item Integrate using trapezoidal rule on the sphere surface:
    \begin{equation}
    \mathcal{C}(r) = \frac{1}{2\pi} \sum_{i,j} \Omega_r(\theta_i, \phi_j) \cdot r^2 \sin\theta_i \, \Delta\theta \, \Delta\phi
    \end{equation}
\end{enumerate}

\textbf{Convergence test:} Verified for radii up to $r = 0.10$ \AA$^{-1}$ (within distance to nearest other crossing)

\textbf{Result:} $\mathcal{C} = +1$ to numerical precision for all tested radii
\end{derivation}

\subsection{Nielsen-Ninomiya Constraint}

\begin{keyresult}[title={Partner Weyl Points}]
The Nielsen-Ninomiya theorem requires:
\begin{equation}
\sum_{\text{all Weyl points}} \chi_i = 0
\end{equation}

\textbf{Our finding:}
\begin{itemize}
    \item 4 Weyl points along $\langle 111 \rangle$ directions with $\chi = +1$
    \item 4 partner Weyl points near $W$ points with $\chi = -1$
    \item Total chirality: $4(+1) + 4(-1) = 0$ $\checkmark$
\end{itemize}

Located by $40 \times 40 \times 40$ grid scan over full BZ.
\end{keyresult}

%% ============================================================
\section{$k \cdot p$ Theory for Weyl Point}
%% ============================================================

\subsection{Two-Band Effective Hamiltonian}

Near the Weyl point, project onto bands 2 and 3:

\begin{equation}
\boxed{
H_{k\cdot p}(\delta\bm{k}) = \omega_W \mathbb{1}_2 + v_\parallel \delta k_\parallel \sigma_z + v_\perp (\delta k_x \sigma_x + \delta k_y \sigma_y)
}
\end{equation}

\begin{derivation}[title={Symmetry Constraints}]
Why this form?

\textbf{$C_3$ symmetry along $[111]$:}
\begin{itemize}
    \item $\sigma_z$ term is $C_3$-invariant (diagonal in the basis)
    \item $\sigma_{x,y}$ terms form the 2D irrep of $C_3$
    \item Bands 2 and 3 belong to \emph{different} irreps ($A$ and $E$)
    \item Therefore $\sigma_{x,y}$ terms cannot mix them at $k_\parallel = 0$
    \item The crossing is \emph{protected} --- no mass term $\propto \sigma_z$ can arise from $C_3$-preserving perturbations
\end{itemize}
\end{derivation}

\subsection{Chirality from $k \cdot p$}

\begin{equation}
\chi = \text{sgn}(v_\parallel v_\perp^2) = \text{sgn}(v_\parallel)
\end{equation}

For our case: $v_\parallel > 0$ so $\chi = +1$, consistent with Berry integration.

\subsection{Fitting Procedure}

\begin{derivation}[title={$k \cdot p$ Parameter Extraction}]
\textbf{Method:}
\begin{enumerate}
    \item Extract LSWT band energies in a window of $\pm 8\%$ of $|\Gamma L|$ centered on Weyl point
    \item Fit two-band eigenvalues $E_\pm = \omega_W \pm \sqrt{v_\parallel^2 \delta k_\parallel^2 + v_\perp^2 \delta k_\perp^2}$
    \item Minimize RMS deviation using Nelder-Mead simplex
\end{enumerate}

\textbf{Results at $D/J = 0.32$:}
\begin{itemize}
    \item $t_W \approx 0.55$ (fractional position along $\Gamma \to L$)
    \item $\omega_W \approx 29$ meV
    \item $v_\parallel \approx v_\perp \approx 10$ meV$\cdot$\AA\ (nearly isotropic cone)
\end{itemize}

\textbf{Agreement:} Direct band-structure fits and $k \cdot p$ fits agree within $\sim 10\%$
\end{derivation}

%% ============================================================
\section{Surface Arc Calculation}
%% ============================================================

\subsection{Slab Geometry}

\begin{derivation}[title={Slab Construction}]
\textbf{Setup:}
\begin{itemize}
    \item 50 unit cells stacked along $a_3$ ([110] direction)
    \item Open boundary conditions at top/bottom surfaces
    \item Periodic boundary conditions in the two in-plane directions
\end{itemize}

\textbf{Hamiltonian:}
\begin{itemize}
    \item Fourier transform only in 2 periodic directions
    \item Retain real-space hopping along finite direction
    \item Results in $200 \times 200$ matrix at each $\bm{k}_\parallel$
\end{itemize}

\textbf{Surface spectral weight:}
\begin{equation}
w_{\text{surf}}(n, \bm{k}_\parallel) = \sum_{l \in \text{surface}} |\psi_n(l, \bm{k}_\parallel)|^2
\end{equation}
where sum is over the outermost 2 unit cells on each surface.
\end{derivation}

\begin{keyresult}[title={Surface Arc Detection}]
Surface arcs are visible as branches with elevated surface weight that:
\begin{itemize}
    \item Connect projections of Weyl points onto surface BZ
    \item Cross the bulk Weyl energy $\omega_W \approx 29$ meV
    \item Are topologically protected (cannot be removed without closing bulk gap)
\end{itemize}
\end{keyresult}

%% ============================================================
\section{Magnon Hall Conductivity}
%% ============================================================

\subsection{Kubo Formula for Thermal Hall}

\begin{equation}
\boxed{
\kappa^{xy} = -\frac{k_B^2 T}{\hbar} \frac{1}{V} \sum_{n,\bm{k}} c_2(n_B(\varepsilon_{n\bm{k}})) \, \Omega_{n\bm{k}}^{xy}
}
\end{equation}

\begin{derivation}[title={Component Definitions}]
\textbf{Bose-Einstein distribution:}
\begin{equation}
n_B(\varepsilon) = \frac{1}{e^{\varepsilon/k_B T} - 1}
\end{equation}

\textbf{Magnon Hall weight function:}
\begin{equation}
c_2(x) = (1+x)\left[\ln\frac{1+x}{x}\right]^2 - (\ln x)^2 - 2\,\text{Li}_2(-x)
\end{equation}
where $\text{Li}_2$ is the polylogarithm.

This function weights the Berry curvature contribution by the thermal occupation.
\end{derivation}

\subsection{Comparison with Experiment}

\begin{assumption}[title={Limitations of LSWT Magnon Hall Calculation}]
\textbf{Observed discrepancy:} Theory exceeds experiment by factor $\sim 200$

\textbf{Reasons for overestimate:}
\begin{enumerate}
    \item Magnon-phonon scattering (not in LSWT) reduces $|\kappa^{xy}|$
    \item Magnon-magnon interactions absent
    \item Finite magnon lifetime from extrinsic scattering
    \item Comparison is with Lu$_2$V$_2$O$_7$, not Y$_2$V$_2$O$_7$
\end{enumerate}

\textbf{Qualitative agreement:}
\begin{itemize}
    \item Both theory and experiment show rise from zero at low $T$
    \item Confirms Berry curvature from Weyl magnon as dominant source
\end{itemize}
\end{assumption}

%% ============================================================
\section{RIXS Formalism}
%% ============================================================

\subsection{Kramers-Heisenberg Formula}

\begin{equation}
\boxed{
\frac{d^2\sigma}{d\Omega\, d\omega_f} = \sum_f \left|
\sum_\nu \frac{\langle f | \mathcal{O}_{\epsilon_f}^\dagger | \nu \rangle \langle \nu | \mathcal{O}_{\epsilon_i} | g \rangle}
{E_g + \hbar\omega_i - E_\nu + i\Gamma_c/2}
\right|^2 \delta(E_f - E_g - \hbar\omega)
}
\end{equation}

\begin{derivation}[title={Physical Interpretation}]
\textbf{States:}
\begin{itemize}
    \item $|g\rangle$: Ground state (energy $E_g$)
    \item $|\nu\rangle$: Intermediate state with core hole (energy $E_\nu$)
    \item $|f\rangle$: Final state (energy $E_f$)
\end{itemize}

\textbf{Process at V $L_3$ edge:}
\begin{enumerate}
    \item Incident photon ($\hbar\omega_i \approx 518$ eV) excites $2p_{3/2} \to 3d$
    \item Intermediate state: $2p^5 3d^2$ configuration
    \item Core hole decays: $3d \to 2p$ photon emission
    \item Final state: excited $3d$ configuration
\end{enumerate}

\textbf{Resonance:} Cross-section peaks when $E_g + \hbar\omega_i \approx E_\nu$
\end{derivation}

\subsection{Polarization Selection Rules}

\begin{keyresult}[title={Channel Separation}]
\textbf{$\sigma \to \pi$ (spin-flip):}
\begin{itemize}
    \item Effective operator carries $\Delta m = \pm 1$
    \item Selects $\Delta S_z = \pm 1$ final states
    \item Contains: Exchange peak (8 meV), SOC peak (22 meV)
\end{itemize}

\textbf{$\sigma \to \sigma$ (spin-conserving):}
\begin{itemize}
    \item Symmetric operator, $\Delta S_z = 0$
    \item Contains: Phonon sideband (100 meV)
\end{itemize}

This polarimetric separation is the \textbf{central experimental prediction}.
\end{keyresult}

%% ============================================================
\section{V$^{4+}$--V$^{4+}$ Dimer Exact Diagonalization}
%% ============================================================

\subsection{Hilbert Space Construction}

\begin{derivation}[title={State Counting}]
\textbf{Single-site initial states:}
\begin{equation}
\binom{10}{1} = 10 \text{ states (one electron in 5 $d$ orbitals $\times$ 2 spins)}
\end{equation}

\textbf{Two-site initial states:}
\begin{equation}
10 \times 10 = 100 \text{ states before symmetry}
\end{equation}
After including exchange coupling and SOC: 190 states

\textbf{Intermediate states (one core hole):}
\begin{itemize}
    \item Core hole can be on either site
    \item Configuration: $2p^5 3d^2$ at one site, $3d^1$ at other
    \item Total: 4560 states
\end{itemize}
\end{derivation}

\subsection{Hamiltonian Parameters}

\begin{table}[h]
\centering
\begin{tabular}{lcc}
\toprule
Parameter & Value & Notes \\
\midrule
$10Dq$ & 1.90 eV & Octahedral splitting \\
$\delta_{\text{trig}}$ & 30 meV & Trigonal distortion \\
$\zeta_{3d}$ & 30 meV & 3d spin-orbit coupling \\
$\zeta_{2p}$ & 4.65 eV & 2p spin-orbit coupling \\
$\Gamma_c$ & 165 meV (HWHM) & Core-hole lifetime \\
$F^2_{dd}$ & 8.102 eV & Slater integral (80\% atomic) \\
$F^4_{dd}$ & 5.083 eV & Slater integral (80\% atomic) \\
$J$ & 8.22 meV & Inter-site exchange \\
$U_{dd}$ & 3.0 eV & Hubbard U \\
\bottomrule
\end{tabular}
\caption{EDRIXS calculation parameters}
\end{table}

\begin{assumption}[title={Slater Integral Scaling}]
Atomic Slater integrals from Cowan's code are scaled to 80\% to account for configuration interaction screening. This is standard practice but introduces $\sim 10$--$20\%$ uncertainty in peak positions.
\end{assumption}

\subsection{Phonon Coupling}

\begin{derivation}[title={Displaced Oscillator Model}]
\begin{equation}
\mathcal{H}_{\text{ph}} = \hbar\omega_{\text{ph}} b^\dagger b + g(b + b^\dagger) \hat{n}_d
\end{equation}

\textbf{Parameters:}
\begin{itemize}
    \item $\hbar\omega_{\text{ph}} = 100$ meV (V--O breathing mode)
    \item $g = 35$ meV (linear coupling)
    \item Huang-Rhys factor: $S_{HR} = g^2/(\hbar\omega_{\text{ph}})^2 \approx 0.12$
\end{itemize}

\textbf{Truncation:} Phonon Hilbert space limited to 4 quanta (converged at this coupling)
\end{derivation}

%% ============================================================
\section{Summary of Key Predictions}
%% ============================================================

\begin{table}[h]
\centering
\begin{tabular}{lccc}
\toprule
Observable & Predicted Value & Current Resolution & Detectable? \\
\midrule
Weyl energy $\omega_W$ & 29 meV & 30 meV FWHM & Marginal \\
Exchange peak & 8 meV & 30 meV FWHM & No (need $<10$ meV) \\
SOC peak & 22 meV & 30 meV FWHM & Yes \\
Phonon sideband & 100 meV & 30 meV FWHM & Yes \\
\bottomrule
\end{tabular}
\caption{Experimental detectability at current RIXS resolution}
\end{table}

\begin{keyresult}[title={Main Conclusions}]
\begin{enumerate}
    \item \textbf{Weyl magnon exists} at $\omega_W \approx 29$ meV with $\mathcal{C} = +1$
    \item \textbf{Protected by $C_3$ symmetry} --- robust to parameter uncertainties
    \item \textbf{Polarimetric RIXS} can separate magnetic (spin-flip) from lattice (spin-conserving) excitations
    \item \textbf{SOC and phonon peaks} testable now; exchange peak requires next-gen resolution
\end{enumerate}
\end{keyresult}

%% ============================================================
\section{Potential Questions and Responses}
%% ============================================================

\subsection{On LSWT Validity}

\textbf{Q: How can LSWT be trusted for $S = 1/2$?}

A: While quantum corrections are nominally $O(1)$, the \emph{topological} aspects are protected:
\begin{itemize}
    \item Weyl point existence: protected by $C_3$ symmetry (cannot be gapped)
    \item Chern number: topologically quantized (cannot change without gap closing)
    \item Interactions shift $\omega_W$ and $t_W$ but preserve the crossing
    \item Prior $1/S$ studies show $\sim 20\%$ renormalization, not $100\%$
\end{itemize}

\subsection{On Parameter Transfer}

\textbf{Q: Why use Lu$_2$V$_2$O$_7$ parameters for Y$_2$V$_2$O$_7$?}

A: Best available data. Key mitigations:
\begin{itemize}
    \item Similar $T_C$ suggests similar $J$
    \item DM ratio affects position but not existence of Weyl point
    \item Topology persists over $D/J \in [0.1, 0.6]$ --- robust to $\pm 20\%$ uncertainty
\end{itemize}

\subsection{On Magnon Hall Discrepancy}

\textbf{Q: Why is theory 200$\times$ larger than experiment?}

A: LSWT gives the \emph{intrinsic} contribution only. Missing physics:
\begin{itemize}
    \item Magnon-phonon scattering
    \item Magnon-magnon interactions
    \item Finite magnon lifetime
    \item Comparison is with different compound (Lu vs Y)
\end{itemize}
The qualitative trend (rise from zero, Berry curvature origin) is correct.

\subsection{On RIXS Momentum Resolution}

\textbf{Q: Can RIXS actually probe the Weyl point at finite $k$?}

A: Current soft X-ray RIXS achieves $\sim 0.03$ \AA$^{-1}$ momentum resolution, sufficient to map dispersion across significant BZ fraction. The dimer calculation gives the zone-center response; momentum-resolved extension is left to future work.

\end{document}
